% -*- coding: utf-8 -*-
% 两阶段 ε-容差 + AHP 加权求解流程图（问题一/二共用）
\begin{figure}[htbp]
\centering
\begin{tikzpicture}[
    node distance=0.7cm and 1.2cm,
    box/.style={draw, rounded corners=3pt, minimum width=3.2cm, minimum height=0.8cm,
        align=center, font=\small, fill=blue!8},
    box2/.style={draw, rounded corners=3pt, minimum width=3.2cm, minimum height=0.8cm,
        align=center, font=\small, fill=green!10},
    box3/.style={draw, rounded corners=3pt, minimum width=3.2cm, minimum height=0.8cm,
        align=center, font=\small, fill=orange!12},
    arrow/.style={->, >=stealth, thick, draw=gray!70},
]

% 第一行：路径生成
\node[box] (gen) {Clarke-Wright 节省算法\\分层路径生成};
\node[box, right=of gen] (fuel) {燃油可行性检查\\加油中转点嵌入};
\node[box, right=of fuel] (pool) {候选路径池 $\mathcal{R}$};

% 第二行：阶段一
\node[box2, below=1.0cm of gen] (s1) {\textbf{阶段一}：$\min T=\sum_r t_r x_r$};
\node[box2, right=of s1] (tstar) {得 $T^*$, 锚点\\$(\tau_0, W_0, F_0)$};

% 第三行：阶段二
\node[box3, below=1.0cm of s1] (s2) {\textbf{阶段二}：$\min Z_2$\\$T\le(1+\varepsilon)T^*$};
\node[box3, right=of s2] (ahp) {AHP 权重加权\\$w_\tau, w_u, w_F$};

% 第四行：输出
\node[box, below=1.0cm of s2] (out) {人员分配与结果输出};

% 箭头
\draw[arrow] (gen) -- (fuel);
\draw[arrow] (fuel) -- (pool);
\draw[arrow] (pool.south) |- (tstar.east);
\draw[arrow] (s1) -- (tstar);
\draw[arrow] (tstar.south) |- (0,-3) -| (s2.north);
\draw[arrow] (s1) -- (s2);
\draw[arrow] (s2) -- (ahp);
\draw[arrow] (s2) -- (out);
\draw[arrow] (ahp.south) |- (out.east);

\end{tikzpicture}
\caption{两阶段 $\varepsilon$-容差 + AHP 加权求解流程}
\label{fig:two_stage_flow}
\end{figure}
